% =====================================================================
%  Chapter 13 — The Universe: Sixty-One Funds and Their Live Windows
% =====================================================================

\chapter{The Universe: Sixty-One Funds and Their Live Windows}
\label{ch:universe}

The universe\idx{Universe} is a deliberately small and deliberately unglamorous thing:
sixty-one exchange-traded funds, all of them US-listed, all of them
buyable with a few hundred dollars, none of them leveraged, inverse,
synthetic or exotic. This chapter sets it out, shows what it can and
cannot contain at each point in time, and argues for the restrictions.

\section{The list}

The universe is grouped into five sleeves, and the grouping is the
book's investment taxonomy as well as its allocation taxonomy.

\textbf{United States equity (eleven funds).} A broad S\&P 500 tracker, a
total-market tracker, a second S\&P 500 tracker, a large-cap growth
tracker, a small-cap tracker, an equal-weight tracker, a mid-cap tracker,
and four growth/value pairs from two providers.

\textbf{International equity (six funds).} Developed and emerging markets
from two providers, an all-country developed tracker, and two all-country
trackers with different breadth.

\textbf{Sector equity (eleven funds).} The eleven sector funds of the
Select Sector SPDR family: technology, financials, energy, health care,
industrials, consumer discretionary, consumer staples, utilities,
materials, real estate and communication services.

\textbf{Factor (eleven funds).} Momentum, value, quality, minimum
volatility, size, high dividend, an environmental screen, developed-market
value, and two more.

\textbf{Bonds (fourteen funds).} Treasuries at three maturities, three
short-bill and cash vehicles, an aggregate bond fund\idx{Fund}, four credit funds
across the investment-grade and high-yield divide, inflation-protected
bonds, a municipal fund, and a non-US aggregate fund.

\textbf{Real assets and alternatives (eight funds).} Two gold vehicles, two
real-estate funds, one international real-estate fund, two broad commodity
funds, a managed-futures fund, and a real-estate income fund.

Two tickers that were on the original list, a cryptocurrency product and a
European equity product, do not exist in the vendor's data and were
dropped rather than substituted. The universe is what it is: sixty-one
funds.

\section{What the universe looks like over time}

The live window\idx{Live window} of each fund is not a detail. It is the difference between
a backtest and a simulation of something that could have happened.

The panel covers 407 months, from January 1993 to November 2026. The
number of funds with a live price on the panel changes as follows: one
fund in 1993, by which the panel begins; twelve by 1999; thirty by 2005;
thirty-nine by 2007; fifty-three by 2013; and all sixty-one from 2019.

The live window of every fund is shown in \figref{fig:live}, and the number of funds with a price on the panel by year is the record's own first fact.

This has three consequences that run through the whole book.

The first is that the early record is a record of a \emph{narrower}
portfolio. In 1994 the investable universe was one fund. The book's
strategic allocation is therefore scaled to what exists: a sleeve with no
live fund has its weight redistributed over the sleeves that do, and the
system holds more cash than the policy would suggest because the policy's
diversification is not yet available. This is not a modelling choice; it
is what a real investor in 1994 faced.

The second is that the model's cross-section is small early on. A
cross-sectional regression needs more than one asset to have any
dispersion in the dependent variable, and needs more still to have a
usable one. The book's model requires a minimum number of training rows
and a minimum number of months before it will produce a forecast at all,
and in the early years it produces nothing. The system then holds the
policy.

The third is that the \emph{record of the strategy} and the \emph{record of
the universe} are different objects, and the book reports the live window
of every fund so that a reader can see exactly which strategy was being
run when.

\section{Why the restrictions are restrictions}

Every restriction in the universe is doing work.

\textbf{No individual stocks.} A book about a single investor who might
buy a hundred shares of something is a different book. The advantage of
the ETF universe is that every position is liquid, that every position has
a published expense ratio, that every position is a known factor or a
known sector, and that the diversification is real rather than assumed.
The cost is that the cross-section is small and the signals are
low-resolution.

\textbf{No options, futures, or crypto.} The owner's scope, and a scope
that suits the thesis. The whole argument of Parts V and VI is that the
protection layer works by reducing exposure and holding cash, and it works
that way because there are no derivatives to hedge with. A book that
introduced a hedging overlay would be testing a different system.

\textbf{No leveraged or inverse products.} These would make the
volatility story artificial: a levered fund's realised volatility is a
designed constant times the index's, which would make the governor's
behaviour trivially predictable and the risk estimates meaningless.

\textbf{Several funds tracking the same index.} The universe deliberately
contains near-duplicates --- three S\&P 500 trackers, two total-market
trackers, two gold funds, two broad commodity funds, three Treasury
funds. This is deliberate for a specific reason: the model's job inside a
sleeve is to choose among funds that are close substitutes, and a universe
with only one fund per index would make the tilt meaningless. It also
raises the cost of a naive equal-weight allocation, which is the right
baseline.

\section{The cost of the universe}

The universe's expense ratios are real and are reported in full in
Appendix C. The mean across the sixty-one funds is a little under twenty
basis points a year; the cheapest are at three and the most expensive at
eighty-five. Those are numbers the reader already pays if they hold any
of these funds, and they are already inside the return series.

The universe's diversification is not complete either. There is no
emerging-market fund beyond the two broad ones, no inflation-linked
outside the United States, no private credit, no real estate outside the
two funds, no infrastructure, no gold beyond the two vehicles. A reader
whose circumstances permit more would extend the universe and would have
to re-run the whole pipeline; the book's code is arranged so that
extending the universe is a matter of adding a row to a list and re-running
the fetch, the panel and the system.

\begin{takeaway}[Why sixty-one]
The universe is restricted so that the system measures what it claims to
measure. Within the restriction it is deliberately broad enough for the
tilt to be meaningful, and it contains near-duplicate funds on purpose, so
that the equal-weight comparator is a serious alternative rather than a
straw man.
\end{takeaway}

\begin{keytakeaways}
\item Sixty-one funds in five sleeves: eleven US equity, six international
equity, eleven sector, eleven factor, fourteen bonds, eight real assets
and alternatives.
\item The live cross-section grows from one fund in 1993 to sixty-one in
2019, and the book's allocation is scaled to what exists at each date.
\item The restrictions --- no stocks, no derivatives, no leverage\idx{Leverage} --- are
what make the protection mechanism in Part V a real mechanism rather than
a hedge overlay.
\item Near-duplicate funds are included on purpose, so that the equal-weight
comparator is a genuine alternative.
\end{keytakeaways}
